\chapter{Conclusion}
\label{chap:conclusion}

\section{Summary of Contributions}

This thesis investigated whether the KL 
convexity guarantee (that a mixture over 
a credal set of posteriors performs at 
least as well as the best individual 
component, remains empirically visible 
when posteriors are approximated by 
flow matching models and mixture weights 
are selected by held-out negative log 
likelihood. The answer is yes, under 
the right architectural and statistical 
conditions.

Four concrete contributions were made 
in pursuit of this question. First, a 
shared-backbone flow matching architecture 
conditioned on a discrete prior index 
was designed and validated, producing 
posterior approximations with comparable 
quality across all $K$ priors 
simultaneously and outperforming 
separately trained flows on predictive 
performance, training stability, and 
parameter efficiency. Second, held-out 
NLL was established as the correct 
weight selection objective, replacing 
proxy divergence objectives that 
introduce misalignment between the 
selection criterion and the theoretical 
guarantee. Third, a systematic empirical 
evaluation across two data regimes, two 
real datasets, and three synthetic 
geometries confirmed all three 
hypotheses: the credal mixture achieves 
near-oracle worst-case regret in the 
prior-sensitive regime (0.0005 nats 
under EM, compared to 0.0039 for 
best-single), provides no benefit or 
harm in the prior-insensitive regime, 
and EM is the most reliable weight 
optimisation algorithm. Fourth, the 
conditions under which the guarantee 
is and is not visible were characterised 
concretely: prior sensitivity is 
necessary, a shared backbone is 
necessary, and likelihood-consistent 
weight selection is necessary. When 
any of these conditions is violated, 
the mixture either degrades or provides 
no improvement over single-prior 
inference.

\section{Limitations}

Several limitations of the present 
work should be acknowledged. The 
primary result is demonstrated on a 
single dataset and problem class, 
Bayesian logistic regression on German 
Credit with $n = 75$. While the 
theoretical framework is general, the 
empirical evidence for the guarantee 
visibility claim rests on this setting, 
and replication on other high-dimensional 
Bayesian problems with genuine prior 
sensitivity is needed before broad 
conclusions can be drawn.

The credal set is finite and 
practitioner-specified. The guarantee 
holds only over the $K$ priors that 
are explicitly included; if the true 
data-generating prior lies outside the 
credal set, the mixture provides no 
additional robustness beyond the best 
included component. The choice of 
which priors to include remains an 
open problem that the present framework 
does not address.

The approximation error $\eta$ in 
Theorem~\ref{thm:approx_mixture_ch2} 
was validated empirically as small 
and approximately uniform across priors 
(roughly 0.015 nats relative to HMC), 
but was not bounded theoretically as 
a function of network capacity or 
training data. The guarantee therefore 
rests on an empirical rather than a 
theoretical characterisation of the 
approximation quality.

Finally, Frank-Wolfe degenerates to 
single-component selection under the 
harmonic step size schedule, and 
mirror descent is sensitive to the 
learning rate. A practitioner who 
selects these algorithms without 
careful tuning would not realise the 
mixture benefit. EM is robust to 
these failure modes, but its 
closed-form update is specific to 
the mixture NLL objective and does 
not generalise to other combination 
criteria.

\section{Future Work}

The most natural extension is to 
continuous credal sets, replacing 
the finite discrete set of $K$ priors 
with a parametric family 
$\{\pi_\lambda : \lambda \in \Lambda\}$. 
This would require integrating the 
mixture weight over a continuous 
simplex, which connects directly to 
the generalised Bayesian inference 
framework of Knoblauch et al.\ (2022) 
and the stochastic interpolant 
formulation of Albergo and 
Vanden-Eijnden (2022). The KL 
convexity guarantee extends to 
continuous mixtures under mild 
regularity conditions, and the 
flow matching architecture is 
well-suited to conditioning on a 
continuous $\lambda$ rather than a 
discrete index.

Preliminary experiments on the Breast 
Cancer Wisconsin Diagnostic dataset 
(UCI, $D = 31$, $n = 75$) provide 
initial evidence that the H1 result 
generalises to higher-dimensional 
medical classification settings: all 
principled mixture algorithms 
outperformed the val-selected 
best-single strategy on all five 
seeds, with mean improvements of 
0.003--0.004 nats. However, in this 
more diffuse posterior regime 
(2.4:1 data-to-parameter ratio) 
uniform weighting outperformed 
optimised weight selection on 
average, suggesting that the 
relative performance of weight 
optimisation algorithms depends on 
the degree of posterior 
concentration. Characterising 
exactly when optimised weighting 
dominates uniform averaging, and 
deriving a criterion for selecting 
between them, is a direction for 
future work.

A second direction is a theoretical 
bound on the approximation gap $\eta$ 
as a function of network capacity, 
training steps, and the divergence 
between the prior and the posterior. 
Such a bound would make the conditions 
in Theorem~\ref{thm:approx_mixture_ch2} 
verifiable without requiring a separate 
MCMC reference, which is the main 
computational overhead of the current 
evaluation pipeline.

Third, the approach could be extended 
to Bayesian neural networks, where 
prior sensitivity has been shown to 
persist even at scale (Fortuin et 
al., 2021; Wenzel et al., 2020). The 
shared backbone architecture scales 
naturally to higher-dimensional 
parameter spaces, and the combination 
of flow matching with credal mixtures 
could provide a principled alternative 
to cold posteriors and other ad hoc 
robustness adjustments currently 
used in deep Bayesian models.

Finally, online credal mixture 
updating, in which mixture weights 
are revised as new data arrives 
rather than computed once on a 
fixed validation set, would extend 
the framework to streaming and 
continual learning settings where 
the most informative prior may 
shift over time.